\documentclass[11pt]{book}
\usepackage[paperwidth=145mm,paperheight=200mm,top=18mm,bottom=20mm,inner=20mm,outer=16mm]{geometry}
\usepackage{brumfiel}
% figures11.tex -- Chapter 11 figures, built 2026-08-17 from the iPad scans.
%
% Every figure below was measured off a scan crop (scan13 pp.19-23, scan14
% pp.3,4,5,7) rather than eyeballed from the photo PDF.  Constrained points
% are DERIVED, never placed by hand:
%   regular polygons -> polar coordinates (k*360/n), never by eye
%   perpendicular bisectors -> (intersection of ...) of two rotated midlines
%   midpoints / apothem feet -> ($(X)!0.5!(Y)$)
%   tangency  -> radius drawn to the point, line laid perpendicular to it
% so the hypothesis stated in the book holds exactly in the drawing.
% Hypotheses are re-asserted numerically in tools/constraints/ch11.py.
%
% Side counts were measured, not guessed.  Fig 11-5's polygon is a HEXAGON:
% fitting the drawn circumscribed and inscribed circles from one common centre
% gives a/R = 0.8646 against cos(pi/6) = 0.8660 (n = 5.97); a heptagon would
% need 0.9010.  See chapters/ch11-UNCERTAIN.md.
%
% Macro names are chapter-prefixed: \FIGXI<NUMBERWORD>.
%
% Line treatment follows the book's own convention: solid where the book
% prints solid (given), dashed where the book prints dashed (aux).

\usetikzlibrary{decorations.pathreplacing,patterns,arrows.meta}

% Primed labels.  The math prime sits ~4pt above the letter, which makes the
% collision audit read it as a separate word overlapping its own letter; a
% 1.3pt drop keeps the pair in one word and is visually indistinguishable.
\newcommand{\XIpr}{\raisebox{-1.3pt}{\ensuremath{{}^{\prime}}}}

% Chapter-local label clearance, as in ch09: a labelled point lying ON a line
% needs a little more outer sep than brumfiel.sty's 3.4pt or the rule clips
% the corner of the label box.
\tikzset{
  vlab/.append style={outer sep=4.8pt},
  slab/.append style={outer sep=4.8pt},
  klab/.append style={outer sep=4.8pt},
  bkarr/.style={fig, -{Stealth[length=3.6pt,width=3pt]}},
  bkarrb/.style={fig, {Stealth[length=3.6pt,width=3pt]}-{Stealth[length=3.6pt,width=3pt]}},
  % Fig 11-6's four callouts are drawn in the book as ARROWS from the word to
  % the part of the polygon it names, so the leader carries a head at the
  % geometry end (the end of the path, since every leader starts at its node).
  lead/.style={fig, line width=0.4pt, -{Stealth[length=3.2pt,width=2.6pt]}},
}

% ------------------------------------------------------- Figs 11-1 and 11-2
% 11-1: circumcentre of a triangle.  m and n are the perpendicular bisectors
% of AB and BC; O is their intersection, hence equidistant from A, B and C.
% AB is drawn horizontal, so m is vertical, exactly as the book prints it.
% 11-2: regular hexagon inscribed in a circle, vertices at k*60 deg with
% C at 0 deg.  Angles 1,2 sit at B and 3,4 at C, each on the bisector of its
% own sub-angle (the 0.33 construction reproduces the book's placement).
\newcommand{\FIGXIONETWO}{\begin{bkfigure}[\linewidth]
\begin{minipage}{0.52\linewidth}\centering
\begin{tikzpicture}[scale=0.86]
  \coordinate (A) at (0,0);
  \coordinate (B) at (2.20,0);
  \coordinate (C) at (4.45,2.00);
  \coordinate (Mab) at ($(A)!0.5!(B)$);
  \coordinate (Mbc) at ($(B)!0.5!(C)$);
  % a second point on each perpendicular bisector, by 90-degree rotation
  \coordinate (mP) at ($(Mab)!1!90:(B)$);
  \coordinate (nP) at ($(Mbc)!1!90:(C)$);
  \coordinate (O)  at (intersection of Mab--mP and Mbc--nP);
  \draw[given] (A) -- (B) -- (C) -- cycle;
  \draw[aux] (O) -- ($(Mab)!-0.10!(O)$);
  \draw[aux] (O) -- ($(Mbc)!-0.18!(O)$);
  \node[vlab, below left] at (A) {$A$};
  \node[vlab, below] at (B) {$B$};
  \node[vlab, above right] at (C) {$C$};
  \node[vlab, above] at (O) {$O$};
  \node[vlab, below] at ($(Mab)!-0.10!(O)$) {$m$};
  \node[vlab, below right] at ($(Mbc)!-0.18!(O)$) {$n$};
\end{tikzpicture}
\figcap{11--1}
\end{minipage}\hfill
\begin{minipage}{0.44\linewidth}\centering
\begin{tikzpicture}[scale=0.86]
  \def\R{1.55}
  \coordinate (O) at (0,0);
  \coordinate (C) at (0:\R);   \coordinate (D) at (60:\R);
  \coordinate (E) at (120:\R); \coordinate (F) at (180:\R);
  \coordinate (A) at (240:\R); \coordinate (B) at (300:\R);
  \draw[given] (O) circle (\R);
  \draw[given] (C) -- (D) -- (E) -- (F) -- (A) -- (B) -- cycle;
  \foreach \p in {A,B,C,D,E} {\draw[given] (O) -- (\p);}
  \fill (O) circle (1.3pt);
  \node[vlab, above left]  at (E) {$E$};
  \node[vlab, above right] at (D) {$D$};
  \node[vlab, right]       at (C) {$C$};
  \node[vlab, below left]  at (A) {$A$};
  \node[vlab, below right] at (B) {$B$};
  \node[vlab, left]        at (O) {$O$};
  \node[slab] at ($($(B)!0.33!(A)$)!0.5!($(B)!0.33!(O)$)$) {1};
  \node[slab] at ($($(B)!0.33!(O)$)!0.5!($(B)!0.33!(C)$)$) {2};
  \node[slab] at ($($(C)!0.33!(B)$)!0.5!($(C)!0.33!(O)$)$) {3};
  \node[slab] at ($($(C)!0.33!(O)$)!0.5!($(C)!0.33!(D)$)$) {4};
\end{tikzpicture}
\figcap{11--2}
\end{minipage}
\end{bkfigure}}

% ------------------------------------------------------- Figs 11-3 and 11-4
% 11-3: l is laid perpendicular to the radius OP at P (direction = OP angle
% + 90), so the tangency hypothesis holds exactly.  A is any other point of l.
% 11-4: l is NOT perpendicular to OP.  m is the perpendicular from O to l,
% A its foot, and B the reflection of P in A -- so AB = AP and OB = OP, which
% is what puts B on the circle.
\newcommand{\FIGXITHREEFOUR}{\begin{bkfigure}[\linewidth]
\begin{minipage}{0.46\linewidth}\centering
% Both halves at 1.30, not 0.86: in 11-4 the letter A has to sit in the wedge
% between OP and l, and at the original scale that wedge was narrower than the
% 9pt letter box.  Same scale on 11-3 so the pair still matches in size.
\begin{tikzpicture}[scale=1.30]
  \def\R{1.25}
  \coordinate (O) at (0,0);
  \coordinate (P) at (-19:\R);
  \coordinate (A) at ($(P)+(71:0.92)$);
  \draw[given] (O) circle (\R);
  \draw[given] (O) -- (P);
  \draw[given] ($(P)+(71:-0.50)$) -- ($(P)+(71:1.84)$);
  \fill (O) circle (1.2pt);
  \fill (A) circle (1.2pt);
  \fill (P) circle (1.2pt);
  \node[vlab, left]        at (O) {$O$};
  \node[vlab, below right] at (P) {$P$};
  \node[vlab, right]       at (A) {$A$};
  \node[vlab, above]       at ($(P)+(71:1.84)$) {$l$};
\end{tikzpicture}
\figcap{11--3}
\end{minipage}\hfill
\begin{minipage}{0.50\linewidth}\centering
\begin{tikzpicture}[scale=1.30]
  \def\R{1.25}
  \coordinate (O) at (0,0);
  \coordinate (A) at (0,-0.39);
  \coordinate (P) at (1.188,-0.39);
  \coordinate (B) at (-1.188,-0.39);
  \draw[given] (O) circle (\R);
  \draw[aux]   (0,1.58) -- (0,-1.58);
  \draw[given] (-1.72,-0.39) -- (1.68,-0.39);
  \draw[given] (O) -- (P);
  \fill (O) circle (1.2pt); \fill (A) circle (1.2pt);
  \fill (B) circle (1.2pt); \fill (P) circle (1.2pt);
  \node[vlab, right]       at (0,0.80) {$m$};
  \node[vlab, above right] at (O) {$O$};
  % A sits in the narrow wedge between OP and l, exactly as the book sets it
  % (scan13 p20).  Anchored placement puts the letter on one line or the
  % other, so the offset is explicit and measured: 0.14 R right, 0.14 R up.
  \node[vlab]              at ($(A)+(0.162,0.191)$) {$A$};
  \node[vlab, above left]  at (B) {$B$};
  \node[vlab, below right] at (P) {$P$};
  \node[vlab, right]       at (1.68,-0.39) {$l$};
\end{tikzpicture}
\figcap{11--4}
\end{minipage}
\end{bkfigure}}

% ------------------------------------------------------- Figs 11-5 and 11-6
% 11-5: regular hexagon (measured, see header) with both circles.  C and E are
% the midpoints of the adjacent sides AB and BD, so OC and OE are apothems and
% the inner circle -- radius R cos(30) -- is tangent at exactly those points.
% 11-6: the four named parts of a regular polygon, with leader lines.
\newcommand{\FIGXIFIVESIX}{\begin{bkfigure}[\linewidth]
\begin{minipage}{0.44\linewidth}\centering
% Both halves are drawn at 1.15 rather than 0.86.  The two figures carry a lot
% of lettering inside the outline (C, E here; four whole words in 11-6), and
% node text does not scale with the picture, so at 0.86 the letters were
% bigger relative to the drawing than the book sets them.
\begin{tikzpicture}[scale=1.15]
  \def\R{1.45}
  \coordinate (O) at (0,0);
  \coordinate (A)  at (84:\R);   \coordinate (B)  at (24:\R);
  \coordinate (D)  at (-36:\R);  \coordinate (V4) at (-96:\R);
  \coordinate (V5) at (-156:\R); \coordinate (V6) at (144:\R);
  \coordinate (C) at ($(A)!0.5!(B)$);
  \coordinate (E) at ($(B)!0.5!(D)$);
  \draw[given] (O) circle (\R);
  \draw[given] (A) -- (B) -- (D) -- (V4) -- (V5) -- (V6) -- cycle;
  \draw[given] (O) circle ({\R*cos(30)});
  \draw[aux]   (O) -- (A);
  \draw[given] (O) -- (C);
  \draw[aux]   (O) -- (B);
  \draw[given] (O) -- (E);
  \draw[aux]   (O) -- (D);
  \fill (O) circle (1.2pt);
  \node[vlab, above]       at (A) {$A$};
  \node[vlab, above right] at (B) {$B$};
  % C and E are tangency points: they sit ON the inscribed circle, so an
  % anchored label grazes that arc.  The book pulls both letters INSIDE the
  % inner circle instead -- C straight down from its side's midpoint, E down
  % and to the left of its own -- measured off scan13 p21 and reproduced here.
  \node[vlab]              at ($(C)+(0.04,-0.34)$) {$C$};
  \node[vlab, left]        at (O) {$O$};
  \node[vlab]              at ($(E)+(-0.29,-0.21)$) {$E$};
  \node[vlab, below right] at (D) {$D$};
\end{tikzpicture}
\figcap{11--5}
\end{minipage}\hfill
\begin{minipage}{0.52\linewidth}\centering
\begin{tikzpicture}[scale=1.15]
  \def\R{1.30}
  \coordinate (O) at (0,0);
  \coordinate (Vr) at (0:\R);
  \coordinate (Vb) at (60:\R);
  \coordinate (Va) at (120:\R);
  \coordinate (Ap) at (-30:{\R*cos(30)});
  \draw[given] (0:\R) -- (60:\R) -- (120:\R) -- (180:\R)
            -- (240:\R) -- (300:\R) -- cycle;
  \draw[aux] (O) -- (Va);
  \draw[aux] (O) -- (Vb);
  \draw[aux] (O) -- (Vr);
  \draw[aux] (O) -- (Ap);
  % The book marks the TOP central angle -- the one between the radii to the
  % 60- and 120-degree vertices -- not the one on the radius side.  Arc radius
  % 0.40 R, measured off scan13 p21.
  \draw[bkarrb] (60:{\R*0.40}) arc (60:120:{\R*0.40});
  \draw (O) circle (0.07);
  \node[slab, align=left, anchor=west] (cang) at (1.30,1.05) {Central\\angle};
  \node[slab, anchor=west] (crad) at (1.46,0.18) {Radius};
  % "Center" sits INSIDE the hexagon at its lower left in the book, with the
  % arrow running up and to the right to the centre mark.
  \node[slab, anchor=east] (ccen) at (0.09,-0.45) {Center};
  \node[slab, anchor=west] (capo) at (1.16,-1.16) {Apothem};
  \draw[lead] (cang.west) -- (66:{\R*0.44});
  \draw[lead] (crad.west) -- ($(O)!0.62!(Vr)$);
  \draw[lead] (-0.62,-0.26) -- (-0.10,-0.115);
  \draw[lead] (capo.west) -- ($(O)!0.72!(Ap)$);
\end{tikzpicture}
\figcap{11--6}
\end{minipage}
\end{bkfigure}}

% ---------------------------------------------------------------- Fig 11-7
% Exterior angle of a regular polygon.  The side DC is produced beyond C; the
% marked angle is between that extension and the next side CB.
\newcommand{\FIGXISEVEN}{\begin{bkfigure}[0.92\linewidth]
\begin{tikzpicture}[scale=0.80]
  \def\R{1.35}
  \coordinate (Cv) at (0:\R);
  \coordinate (Bv) at (60:\R);
  \coordinate (Dv) at (-60:\R);
  \draw[given] (0:\R) -- (60:\R) -- (120:\R) -- (180:\R)
            -- (240:\R) -- (300:\R) -- cycle;
  % Measured off scan13 p22: the extension runs a full side length beyond C,
  % and the marked angle's arc has radius ~0.53 R, not the 0.31 R first drawn.
  \coordinate (Ext) at ($(Cv)!-1.00!(Dv)$);
  \draw[aux] (Cv) -- (Ext);
  \draw[bkarrb] ($(Cv)+(120:0.72)$) arc (120:60:0.72);
\end{tikzpicture}
\figcap{11--7}
\end{bkfigure}}

% ------------------------------------------------------- Figs 11-8 and 11-9
% Two regular hexagons of different size, flat-topped (vertices at k*60 deg
% with C at 0 deg).  11-9 adds Q, the midpoint of the top side AB, so that
% OQ is the apothem a and OA the radius r.
\newcommand{\FIGXIEIGHTNINE}{\begin{bkfigure}[\linewidth]
\begin{minipage}{0.44\linewidth}\centering
% 0.70 here against 0.90 in 11-9: the two blocks share one text line, and 11-9
% needs the larger scale to keep r', a' off the radii it sets them between.
\begin{tikzpicture}[scale=0.70]
  \def\Ra{1.55}\def\Rb{1.10}
  \begin{scope}
    \coordinate (O) at (0,0);
    \coordinate (A) at (120:\Ra); \coordinate (B) at (60:\Ra);
    \coordinate (C) at (0:\Ra);   \coordinate (D) at (-60:\Ra);
    \coordinate (E) at (-120:\Ra);
    \draw[given] (0:\Ra) -- (60:\Ra) -- (120:\Ra) -- (180:\Ra)
              -- (240:\Ra) -- (300:\Ra) -- cycle;
    \draw[aux] (O) -- (A);
    \draw[aux] (O) -- (B);
    \fill (O) circle (1.2pt);
    \node[vlab, above left]  at (A) {$A$};
    \node[vlab, above right] at (B) {$B$};
    \node[vlab, right]       at (C) {$C$};
    \node[vlab, below right] at (D) {$D$};
    \node[vlab, below left]  at (E) {$E$};
    \node[vlab, below]       at (O) {$O$};
  \end{scope}
  \begin{scope}[xshift=3.35cm]
    \coordinate (Op) at (0,0);
    \coordinate (Ap) at (120:\Rb); \coordinate (Bp) at (60:\Rb);
    \coordinate (Cp) at (0:\Rb);   \coordinate (Dp) at (-60:\Rb);
    \coordinate (Ep) at (-120:\Rb);
    \draw[given] (0:\Rb) -- (60:\Rb) -- (120:\Rb) -- (180:\Rb)
              -- (240:\Rb) -- (300:\Rb) -- cycle;
    \draw[aux] (Op) -- (Ap);
    \draw[aux] (Op) -- (Bp);
    \fill (Op) circle (1.2pt);
    \node[vlab, above left]  at (Ap) {$A\XIpr$};
    \node[vlab, above right] at (Bp) {$B\XIpr$};
    \node[vlab, right]       at (Cp) {$C\XIpr$};
    \node[vlab, below right] at (Dp) {$D\XIpr$};
    \node[vlab, below left]  at (Ep) {$E\XIpr$};
    \node[vlab, below]       at (Op) {$O\XIpr$};
  \end{scope}
\end{tikzpicture}
\figcap{11--8}
\end{minipage}\hfill
\begin{minipage}{0.52\linewidth}\centering
% 0.90, not 0.80.  r, a, r' and a' are set inside the wedge between two
% radii, and on the smaller polygon that wedge is only a few points wider
% than the letters themselves; the extra 12% is what keeps them off the lines.
\begin{tikzpicture}[scale=0.90]
  \def\Ra{1.45}\def\Rb{1.02}
  \begin{scope}
    \coordinate (O) at (0,0);
    \coordinate (A) at (120:\Ra); \coordinate (B) at (60:\Ra);
    \coordinate (C) at (0:\Ra);
    \coordinate (Q) at ($(A)!0.5!(B)$);
    \draw[given] (0:\Ra) -- (60:\Ra) -- (120:\Ra) -- (180:\Ra)
              -- (240:\Ra) -- (300:\Ra) -- cycle;
    \draw[aux, bkarr] (A) -- (O);
    \draw[aux, bkarr] (Q) -- (O);
    \draw[aux, bkarr] (B) -- (O);
    \fill (O) circle (1.2pt);
    \node[vlab, above left]  at (A) {$A$};
    \node[vlab, above]       at (Q) {$Q$};
    \node[vlab, above right] at (B) {$B$};
    \node[vlab, right]       at (C) {$C$};
    \node[vlab, below]       at (O) {$O$};
    \node[slab, left]  at ($(A)!0.45!(O)$) {$r$};
    % The wedge between OQ and OB tapers toward O, so a sits high in it, as
    % the book has it (scan13 p23); coordinates solved for equal clearance.
    \node[slab] at (0.186,0.944) {$a$};
  \end{scope}
  \begin{scope}[xshift=3.05cm]
    \coordinate (Op) at (0,0);
    \coordinate (Ap) at (120:\Rb); \coordinate (Bp) at (60:\Rb);
    \coordinate (Cp) at (0:\Rb);
    \coordinate (Qp) at ($(Ap)!0.5!(Bp)$);
    \draw[given] (0:\Rb) -- (60:\Rb) -- (120:\Rb) -- (180:\Rb)
              -- (240:\Rb) -- (300:\Rb) -- cycle;
    \draw[aux, bkarr] (Ap) -- (Op);
    \draw[aux, bkarr] (Qp) -- (Op);
    \draw[aux, bkarr] (Bp) -- (Op);
    \fill (Op) circle (1.2pt);
    \node[vlab, above left]  at (Ap) {$A\XIpr$};
    \node[vlab, above]       at (Qp) {$Q\XIpr$};
    \node[vlab, above right] at (Bp) {$B\XIpr$};
    \node[vlab, right]       at (Cp) {$C\XIpr$};
    \node[vlab, below]       at (Op) {$O\XIpr$};
    % On the smaller hexagon the corridors these two letters live in -- r'
    % between side A'(180-deg vertex) and radius O'A', a' between O'Q' and
    % O'B' -- are only a couple of points wider than the primed letter boxes
    % themselves (5.4 x 10.0 pt, measured off the compiled page).  Anchored
    % placement cannot fit them, so both are set at coordinates solved for
    % equal clearance on every side; r' still hugs the radius it names, as
    % the book has it.
    \node[slab] at (-0.508,0.333) {$r\XIpr$};
    \node[slab] at ( 0.141,0.623) {$a\XIpr$};
  \end{scope}
\end{tikzpicture}
\figcap{11--9}
\end{minipage}
\end{bkfigure}}

% --------------------------------------------------------------- Fig 11-10
% Inscribed square; bisecting the central angle in the first quadrant gives
% the octagon vertex at 45 deg, and the two dashed chords are the octagon
% sides there.  The bisection is exact (45 = 90/2), not drawn by eye.
\newcommand{\FIGXITEN}{\begin{bkfigure}[0.56\linewidth]
\begin{tikzpicture}[scale=0.92]
  \def\R{1.50}
  \coordinate (O) at (0,0);
  \draw[given] (O) circle (\R);
  \draw[given] (90:\R) -- (0:\R) -- (-90:\R) -- (180:\R) -- cycle;
  % Checked against scan14 p3 at 3x: the radii to the square's own vertices
  % carry NO arrowheads; only the bisecting ray does, and that one is
  % double-headed (a head at the centre and a head at the circle).
  \draw[aux] (O) -- (90:\R);
  \draw[aux] (O) -- (0:\R);
  \draw[aux, bkarrb] (O) -- (45:\R);
  \draw[aux] (90:\R) -- (45:\R) -- (0:\R);
\end{tikzpicture}
\figcap{11--10}
\end{bkfigure}}

% --------------------------------------------------------------- Fig 11-11
% Circle of radius 1 with the inscribed hexagon (solid, vertices at 30+k*60)
% and the circumscribed hexagon (dashed, vertices at k*60, circumradius
% R/cos 30) -- so the dashed hexagon's sides are tangent to the circle at the
% solid hexagon's vertices, exactly as the book draws it.
\newcommand{\FIGXIELEVEN}{\begin{bkfigure}[0.52\linewidth]
\begin{tikzpicture}[scale=1.05]
  \def\R{1.30}
  \coordinate (O) at (0,0);
  \draw[given] (O) circle (\R);
  \draw[given] (30:\R) -- (90:\R) -- (150:\R) -- (210:\R)
            -- (270:\R) -- (330:\R) -- cycle;
  \draw[aux] (0:{\R/cos(30)})   -- (60:{\R/cos(30)})  -- (120:{\R/cos(30)})
          -- (180:{\R/cos(30)}) -- (240:{\R/cos(30)}) -- (300:{\R/cos(30)})
          -- cycle;
  \draw[bkarr] (O) -- (0:\R);
  \fill (O) circle (1.2pt);
  \node[slab, below] at (0:{\R*0.55}) {1};
\end{tikzpicture}
\figcap{11--11}
\end{bkfigure}}

% --------------------------------------------------------------- Fig 11-12
% A and C are ends of a side of the inscribed polygon; B is its midpoint, so
% OB is the apothem a_n and OB is perpendicular to the chord AC.
\newcommand{\FIGXITWELVE}{\begin{bkfigure}[0.50\linewidth]
\begin{tikzpicture}[scale=1.0]
  \def\R{1.40}
  % A and C at +-45 deg (was +-40): measured off scan14 p5, where the chord
  % subtends ~90 deg and the drawn apothem is 0.69 r, not 0.77 r.
  \coordinate (O) at (0,0);
  \coordinate (A) at (45:\R);
  \coordinate (C) at (-45:\R);
  \coordinate (B) at ($(A)!0.5!(C)$);
  \draw[given] (O) circle (\R);
  \draw[given] (O) -- (A);
  \draw[given] (O) -- (B);
  \draw[given] (O) -- (C);
  \draw[given] (A) -- (C);
  \fill (O) circle (1.2pt);
  \node[vlab, above right] at (A) {$A$};
  % B lies on the chord AC with the circle close behind it: the chapter-wide
  % 4.8pt outer sep pushed the letter onto the arc, so this one label uses the
  % style file's own 3.4pt and sits where the book sets it, between the two.
  \node[vlab, right, outer sep=2.6pt] at (B) {$B$};
  \node[vlab, below right] at (C) {$C$};
  \node[vlab, left]        at (O) {$O$};
  \node[slab, above left]  at ($(O)!0.55!(A)$) {$r$};
  \node[slab, above]       at ($(O)!0.66!(B)$) {$a_n$};
\end{tikzpicture}
\figcap{11--12}
\end{bkfigure}}

% --------------------------------------------------------------- Fig 11-13
% Annular ring.  Radii 2 and 4 as in Exercise 11-13.2, i.e. inner = outer/2.
\newcommand{\FIGXITHIRTEEN}{\begin{bkfigure}[\linewidth]
\begin{tikzpicture}[scale=1.0]
  \def\Ro{1.30}
  \path[pattern=vertical lines, even odd rule]
        (0,0) circle (\Ro) (0,0) circle ({\Ro/2});
  \draw[given] (0,0) circle (\Ro);
  \draw[given] (0,0) circle ({\Ro/2});
\end{tikzpicture}
\figcap{11--13}
\end{bkfigure}}

% --------------------------------------------------------------- Fig 11-14
% ORGANIC SHAPE -- built against scan14 p7 (the straight duplicate), per the
% source policy.  AC is a diameter of the big circle and B its centre, with
% AB = BC.  The S-curve is the semicircle on AB drawn ABOVE the diameter
% followed by the semicircle on BC drawn BELOW it; the shaded region is
% everything inside the big circle below that curve, which is why its area is
% exactly half the disc (Exercise 11-13.3).
\newcommand{\FIGXIFOURTEEN}{\begin{bkfigure}[\linewidth]
\begin{tikzpicture}[scale=0.95]
  \def\R{1.40}
  \coordinate (A) at (-\R,0);
  \coordinate (B) at (0,0);
  \coordinate (C) at (\R,0);
  \path[pattern=north east lines]
        (A) arc[start angle=180, end angle=0,   radius={\R/2}]
            arc[start angle=180, end angle=360, radius={\R/2}]
            arc[start angle=0,   end angle=-180, radius=\R] -- cycle;
  \draw[given] (B) circle (\R);
  \draw[given] (A) arc[start angle=180, end angle=0,   radius={\R/2}];
  \draw[given] (B) arc[start angle=180, end angle=360, radius={\R/2}];
  \fill (A) circle (1.3pt); \fill (B) circle (1.3pt); \fill (C) circle (1.3pt);
  % A stands inside the hatched region; the book clears a small white patch of
  % hatching around the letter rather than printing it over the ruling.
  \node[vlab, right, fill=white, inner sep=1.6pt] at (A) {$A$};
  \node[vlab, right] at (B) {$B$};
  \node[vlab, right] at (C) {$C$};
\end{tikzpicture}
\figcap{11--14}
\end{bkfigure}}

% --------------------------------------------------------------- Fig 11-15
% Square with both its circles: the circumscribed circle through the vertices
% (radius R) and the inscribed circle touching the sides (radius R/sqrt 2).
\newcommand{\FIGXIFIFTEEN}{\begin{bkfigure}[\linewidth]
\begin{tikzpicture}[scale=1.0]
  \def\R{1.30}
  \coordinate (O) at (0,0);
  \draw[given] (O) circle (\R);
  \draw[given] (45:\R) -- (135:\R) -- (225:\R) -- (315:\R) -- cycle;
  \draw[given] (O) circle ({\R/sqrt(2)});
\end{tikzpicture}
\figcap{11--15}
\end{bkfigure}}

% --------------------------------------------------------------- Fig 11-16
% Semicircles erected outward on the three sides of a right triangle; the
% right angle is at the lower-left vertex, so the two legs are the vertical
% and horizontal sides and the hypotenuse carries the largest semicircle.
\newcommand{\FIGXISIXTEEN}{\begin{bkfigure}[\linewidth]
\begin{tikzpicture}[scale=0.86]
  \coordinate (L) at (0,0);
  \coordinate (T) at (0,1.60);
  \coordinate (R) at (2.52,0);
  \coordinate (Mh) at ($(T)!0.5!(R)$);
  \draw[given] (L) -- (T) -- (R) -- cycle;
  \draw[aux] (T) arc[start angle=90,  end angle=270, radius=0.80];
  \draw[aux] (L) arc[start angle=180, end angle=360, radius=1.26];
  \draw[aux] (T) arc[start angle={atan2(0.80,-1.26)},
                     end angle={atan2(0.80,-1.26)-180}, radius=1.4925];
\end{tikzpicture}
\figcap{11--16}
\end{bkfigure}}


% ---- local macros (hoist into book preamble) ----
% Theorem head carrying an unparenthesised note.  The book prints
% "Theorem 11-4 Converse to Theorem 11-3." with no parentheses, but
% brumfiel.sty's bkplain style wraps \thmnote in parentheses.  Shares the
% `theorem' counter so the numbering stays in one sequence.
\newtheoremstyle{XIplainnote}{1.1\topsep}{1.1\topsep}{\itshape}{}%
  {\bfseries}{.}{ }{\thmname{#1}\thmnumber{ #2}\thmnote{ \normalfont#3}}
\theoremstyle{XIplainnote}
\newtheorem{XIthmnote}[theorem]{Theorem}
\theoremstyle{bkplain}
% (No star-convention footnote here: ch11 has starred exercises -- Ex. Gp.
% 11-3 #2, 11-8 ##8-10, 11-12 ##2-3, 11-13 #13 -- but the book prints no such
% footnote on any page of the chapter, book pp. 196, 200, 205, 207.
% ch02:311 owns the book-wide first use; same call as ch10.)
% Numbered remarks. The book sets italic "Remark" + roman number with NO
% period after the word ("Remark 1."), unlike the shared \begin{remark} in
% brumfiel.sty, which prints "Remark." Verified against book pp. 196, 205.
\newenvironment{XIremarkn}[1]{\par\medskip\noindent\emph{Remark} #1.\ }%
  {\par\medskip}
% Plural run-in label, as the book sets the note under Theorem 11-12
% (book p. 201) -- same case as ch14's \XIVremarks.
\newenvironment{XIremarks}{\par\medskip\noindent\emph{Remarks.}\ }%
  {\par\medskip}
% Footnotes marked with the book's symbols (* then \dag) rather than digits;
% both fall on book p. 204.  Same device as ch12's \XIIsymfoot.
\newcommand{\XIsymfoot}[1]{%
  \begingroup\renewcommand{\thefootnote}{\fnsymbol{footnote}}%
  \footnote{#1}\endgroup}
% The book sets "n >= 3" with the double-bar relation, not \geq.
\newcommand{\XIgeqq}{\geqq}
% ---- end local macros ----

\begin{document}
\setcounter{chapter}{11}

\chapter*{\raggedright\Huge CIRCLES AND\\REGULAR POLYGONS}
\markboth{CIRCLES AND REGULAR POLYGONS}{}
\addcontentsline{toc}{chapter}{11.\ Circles and Regular Polygons}
\vspace{-1em}\noindent\rule{\linewidth}{0.6pt}\vspace{1em}
\figurekey

%========================================================= 11-1
\section*{11--1\quad INTRODUCTION}
\addcontentsline{toc}{section}{11--1\quad Introduction}

In this chapter we define what we mean by the \emph{circumference} (or
\emph{length}) of a \emph{circle} and then prove that the circumference $C$ of
a circle of radius $r$ is given by the familiar formula, $C = 2\pi r$. This
involves giving a definition of the number $\pi$. To do these things, we must
learn a few facts about limits, because the study of circumference is made by
approximating to the circle by polygons and considering polygons with a very
large number of sides. It is therefore necessary, first, to prove some theorems
about regular polygons.

%========================================================= 11-2
\section*{11--2\quad CIRCLES AND REGULAR POLYGONS}
\addcontentsline{toc}{section}{11--2\quad Circles and regular polygons}

Our study of the circle depends upon knowing about regular polygons inscribed
in circles. We therefore study these matters first.

\begin{definition}
A regular polygon is \emph{inscribed} in a circle if all the vertices of the
polygon are on the circle. We also say that the circle \emph{circumscribes the
polygon}.
\end{definition}

\exercisegroup{11--1}
\begin{exercisecols}
\begin{exlist}
\item With ruler and compass, inscribe a square in a given circle.
\item Using Exercise 1, show how to construct regular inscribed polygons of
      8 sides; 16 sides; 32 sides.
\item With ruler and compass, inscribe a regular hexagon (6-sided polygon) in
      a given circle.
\item Using Exercise 3, show how to construct regular polygons of 3 sides;
      12 sides; 24 sides; 48 sides.
\end{exlist}
\end{exercisecols}

\begin{theorem}
There is exactly one circle circumscribing a given triangle.
\end{theorem}

\noindent\emph{Proof.} In Fig.~11--1, suppose that the triangle is
$\tri{ABC}$, and choose any two of the sides, say $AB$ and $BC$. Let $m$ and
$n$ be the perpendicular bisectors of these sides. The lines $m$ and $n$ must
intersect at a point $O$, for if $m$ and $n$ did not intersect, they would be
parallel, $CB$ would be perpendicular to $m$, and the points $A$, $B$, and $C$
would be collinear. The point $O$ is equidistant from $A$ and $B$, and also
equidistant from $B$ and $C$. (Why?) Hence $\sg{OA} = \sg{OB} = \sg{OC}$, and
the circle with center $O$ and radius $\sg{OA}$ passes through $A$, $B$, and
$C$. This proves that there is at least one circle through $A$, $B$, and $C$.
To see that there can be only one, we observe that the center of such a circle
must be equidistant from $A$, $B$, and $C$, and hence must lie on the
perpendicular bisectors of $AB$ and $BC$. These lines, $m$ and $n$, intersect
in only one point. This completes the proof.

\FIGXIONETWO

\begin{theorem}
There is exactly one circle that circumscribes a regular polygon.
\end{theorem}

\noindent\emph{Proof.} In Fig.~11--2, suppose the polygon has successive
vertices $A$, $B$, $C$, $D,\ldots$. We show that the circle circumscribing
$\tri{ABC}$ actually circumscribes the polygon. Let $O$ be the center of the
circle circumscribing $\tri{ABC}$, and consider the segments $OA$, $OB$, $OC$,
$OD$, $OE,\ldots$. By definition of circumscribing circle,
$\sg{OA} = \sg{OB} = \sg{OC} = r$. We show that $\sg{OD} = r$. The same method
can be used to prove that $\sg{OE} = \cdots = r$. Hence we need only show that
$\sg{OD} = r$.

\begin{steps}
\step{$\ang{2} \cong \ang{3}$.}{Why?}
\step{$\ang{ABC} \cong \ang{BCD}$, and $AB \cong BC$.}{The polygon is regular.}
\step{$\ang{1} \cong \ang{4}$.}{Addition and subtraction of angles postulate.}
\step{$\tri{OAB} \cong \tri{ODC}$.}{S.A.S.}
\step{$OD \cong OA$.}{Statement 4.}
\step{$\sg{OD} = r$.}{Statement 5.}
\end{steps}

\exercise{11--2}
Suppose that in Theorem 11--1 we had considered the perpendicular bisectors of
$AB$ and $AC$ instead of $AB$ and $BC$. Would the same point $O$ have been
found? What does the theorem say about the perpendicular bisectors of the
sides of a triangle?

%========================================================= 11-3
\section*{11--3\quad TANGENTS TO CIRCLES}
\addcontentsline{toc}{section}{11--3\quad Tangents to circles}

To talk about circles \emph{inscribed} in regular polygons, we need to know
about tangents to circles.

\begin{theorem}
If $OP$ is a radius of a circle with center $O$, and a line $l$ passes through
$P$ and is perpendicular to $OP$, then $l$ has exactly one point in common with
the circle.
\end{theorem}

\noindent\emph{Proof.} In Fig.~11--3, let $A$ be any point on $l$ different
from $P$. Then $\tri{OPA}$ is a right triangle, and $\sg{OA} > \sg{OP}$.
(Why?) Hence $A$ is outside the circle, and $l$ intersects the circle only at
the point $P$.

\begin{definition}
A line $l$ is \emph{tangent} to a circle if it has exactly one point in common
with the circle; the point is called \emph{the point of tangency}.
\end{definition}

\FIGXITHREEFOUR

\begin{XIthmnote}[Converse to Theorem 11--3]
A line tangent to a circle at a point $P$ is perpendicular to the radius
containing $P$.
\end{XIthmnote}

\noindent\emph{Proof.} We will prove the contrapositive, that is, if the line
is not perpendicular to the radius at $P$, then it is not tangent to the
circle. Consider a circle (Fig.~11--4) with center $O$ and a point $P$ on the
circle. Let $l$ be a line through $P$, such that $l$ is not perpendicular to
$OP$.

\begin{steps}
\step{Let $m$ be the line $\perp l$ passing through $O$. Let $A$ be the
      intersection of $m$ and $l$.}{Why is there such a line?}
\step{Let $B$ be the point on $l$ on the opposite side of $A$ from $P$, and
      such that $\sg{AB} = \sg{AP}$.}{Why does such a point exist?}
\step{$\sg{OB} = \sg{OP}$.}{Why? Several different arguments may be given.}
\step{$B$ is on the circle and also on $l$.}{Statements 2 and 3 and definition
      of a circle.}
\step{$l$ intersects the circle twice, and hence $l$ is not tangent to the
      circle. This completes the proof.}{Statement 4 and definition of
      tangent.}
\end{steps}

\begin{XIremarkn}{1}
In the above proof, the point $A$ had to be inside the circle because
$\sg{OA} < \sg{OP}$, by the Pythagorean theorem. The Pythagorean theorem shows
that the perpendicular distance from a point to a line is less than the
distance to any other point on the line.
\end{XIremarkn}

\begin{XIremarkn}{2}
An argument similar to the proof of Theorem 11--4, and using the Pythagorean
theorem, can be made to show that a line can intersect a circle at most twice.
Hence, a line will either be tangent to a circle, or intersect the circle in
exactly two points, or fail to intersect the circle.
\end{XIremarkn}

\exercisegroup{11--3}
\begin{exercisecols}
\begin{exlist}
\item Draw a circle, and construct (\RC) the line tangent to the circle at a
      point of the circle.
\item[\stex 2.] From a point outside a circle, construct a line tangent to the
      circle.
\item[3.] Prove that if $A$ is a point not on a line $l$, and $B$ and $C$ are
      distinct points on $l$, with $AB \perp l$, then $\sg{AB} < \sg{AC}$.
\end{exlist}
\end{exercisecols}

\begin{definition}
A circle is said to be \emph{inscribed} in a polygon, if it is tangent to each
of the sides of the polygon. We also say that the polygon \emph{circumscribes
the circle}.
\end{definition}

\begin{theorem}
There is a circle inscribed in every regular polygon.
\end{theorem}

\begin{steps}
\step{Let $O$ be the center of the circle circumscribing the polygon
      (Fig.~11--5).}{Theorem 11--2.}
\step{Let $AB$ be a side of the polygon, and $C$ the midpoint of $AB$.}{$AB$
      has one, and only one, midpoint.}
\step{$OC \perp AB$.}{Why?}
\step{The circle with center $O$ and radius $OC$ is tangent to $AB$ at
      $C$.}{Theorem 11--4.}
\step{Now let $BD$ be a side of the polygon adjacent to $AB$, and $E$ the
      midpoint of $BD$. Then $OE \cong OC$.}{Why? This takes several steps.}
\step{Hence the circle with radius $OC$ is also tangent to $BD$ at
      $E$.}{Statement 5 and the fact that $OE \perp BD$.}
\end{steps}

\noindent 7. The above argument can be continued around the polygon to complete
the proof.

\FIGXIFIVESIX

\exercise{11--4}
Show that congruent chords in a circle are equidistant from the center of the
circle.

\begin{definition}
The \emph{center} of a regular polygon is the center of the inscribed or
circumscribed circle. The \emph{apothem} of a regular polygon is the length of
the radius of the inscribed circle (that is, a number). The \emph{radius} of a
regular polygon is the radius of the circumscribed circle. A \emph{central
angle} of a regular polygon is an angle, with vertex at the center, whose sides
contain adjacent vertices of the polygon. See Fig.~11--6.
\end{definition}

% multicols dropped: this group carries Fig 11-7, which now sits inline after
% the exercise (no. 4) that names it, instead of being hoisted above the whole
% list.  (Caleb, 2026-08-17: exercise-group figures belong inline.)
\exercisegroup{11--5}
\begin{exercisecols}
\begin{exlist}
\item Inscribe a regular hexagon in a circle. Draw the circle inscribing the
      hexagon.
\item Using your experience in measuring angles in degrees (we will study this
      topic in the next chapter), how many degrees in a central angle of a
      regular pentagon? hexagon? $n$-gon?
\item What is the sum of the degree measures of the angles of a regular
      pentagon? hexagon? $n$-gon?
\item If an exterior angle of a regular polygon is as shown in Fig.~11--7, how
      many degrees in it if the polygon has 5 sides? 6 sides? $n$ sides?
\exfig{\FIGXISEVEN}
\item Prove that a central angle of a regular polygon is the supplement of any
      angle of the polygon.
\item Show that an exterior angle of a regular polygon is congruent to a
      central angle.
\item In Definition 11--4, how do we know that the center of the inscribed
      polygon and the center of the circumscribed polygon are identical?
\end{exlist}
\end{exercisecols}

%========================================================= 11-4
\section*{11--4\quad REGULAR POLYGONS AND SIMILARITY}
\addcontentsline{toc}{section}{11--4\quad Regular polygons and similarity}

All our proofs about area and circumference of circles depend upon the
following theorem.

\begin{theorem}
Two regular polygons are similar if, and only if, they have the same number of
sides.
\end{theorem}

\noindent\emph{Preliminary remark.} Recall that two polygons are similar if
there is a correspondence between their angles and sides, such that
corresponding angles are congruent and corresponding sides are proportional.
Thus, if the polygons are similar, they necessarily have the same number of
sides, and the ``only if'' part of the theorem is obvious. Therefore, we have
only to prove that if they have the same number of sides, they must be similar.

\noindent\emph{Proof.} Suppose that two regular polygons have consecutive
vertices $A$, $B$, $C$, $D$, $E,\ldots$ and $A'$, $B'$, $C'$, $D'$,
$E',\ldots$, respectively, as in Fig.~11--8. Suppose also that the polygons are
not similar. Since they are regular,
\[
\frac{\sg{AB}}{\sg{A'B'}} = \frac{\sg{BC}}{\sg{B'C'}}
= \frac{\sg{CD}}{\sg{C'D'}} = \frac{\sg{DE}}{\sg{D'E'}} = \cdots.
\]
It must be, then, that $\ang{A} \ncong \ang{A'}$, $\ang{B} \ncong \ang{B'}$,
etc. If $\ang{A} \ncong \ang{A'}$, then $\ang{A} > \ang{A'}$ or
$\ang{A} < \ang{A'}$. Suppose that $\ang{A} < \ang{A'}$. Since the polygons are
regular, $\ang{B} < \ang{B'}$, $\ang{C} < \ang{C'}$, etc. The exterior angle at
$A$ is greater than the exterior angle at $A'$. (Why?) This is true also for
the exterior angles at $B$ and $B'$, $C$ and $C'$, etc. The sum of $n$ exterior
angles in polygon $ABCDE\ldots$ is two straight angles. Clearly then, the sum
of $n$ exterior angles in polygon $A'B'C'D'E'\ldots$ is not two straight
angles. This contradiction completes the proof.

\begin{remark}
If we had studied measurement of angles, we could more easily have proved that
$\ang{AOB} \cong \ang{A'O'B'}$, because the measure of these angles depends
only upon the number of sides $n$. In degrees, the measure of $\ang{AOB}$ would
be equal to $360/n$.
\end{remark}

\FIGXIEIGHTNINE

\begin{theorem}
The perimeters of regular polygons, with the same number of sides, are
proportional to their radii (or apothems).
\end{theorem}

\noindent\emph{Proof.} Suppose that two polygons (of $n$ sides) are labeled as
in Fig.~11--9, and have perimeters $p$ and $p'$. We must show that
$p/p' = r/r' = a/a'$.

\begin{steps}
\step{$\tri{AOQ} \simto \tri{A'O'Q'}$.}{Why?}
\step{$r/r' = a/a' = \sg{AQ}/\sg{A'Q'}$.}{Statement 1.}
\step{$p' = n \cdot \sg{A'B'} = 2n \cdot \sg{A'Q'}$, and
      $p = n \cdot \sg{AB} = 2n \cdot \sg{AQ}$.}{Definition of perimeter and
      the fact that $Q$ and $Q'$ are midpoints.}
\step{$p/p' = \sg{AQ}/\sg{A'Q'}$.}{Statement 3 and algebra.}
\step{$p/p' = r/r' = a/a'$.}{Statements 2 and 4.}
\end{steps}

\begin{theorem}
The area of a regular polygon is equal to one-half the product of apothem and
perimeter; area $= \frac{1}{2}a \cdot p$.
\end{theorem}

\noindent The proof is left as an exercise.

\begin{theorem}
The areas of two regular polygons, with the same number of sides, are
proportional to the squares of their radii (or the squares of their apothems,
or the squares of their sides).
\end{theorem}

\noindent The proof is left as an exercise.

\exercisegroup{11--6}
\begin{exercisecols}
\begin{exlist}
\item Two regular heptagons (7-gons) have sides 3 in. and 4 in., respectively.
      What is the ratio of their perimeters? of their areas?
\item A regular pentagon has twice the area of another regular pentagon. What
      is the ratio of their perimeters?
\item Find the radius and apothem of an equilateral triangle of side 2 in.
\item Two regular polygons of the same number of sides have perimeters of 28
      and 42 in. The radius of the first polygon is $x$ in. What is the radius
      of the second?
\item The apothems of two regular hexagons are 3 in. and 5 in. What is the
      ratio of their areas?
\item The area of a regular 17-gon is 16/9 times the area of a second 17-gon.
      What is the ratio of their perimeters? radii?
\item If a regular hexagon and an equilateral triangle are inscribed in a
      circle, the side of the hexagon is twice the apothem of the triangle.
      Show this.
\item Find the areas of the squares inscribed in and circumscribed about a
      circle of radius 1. Also find the perimeters of the squares.
\item Repeat Exercise 8, using hexagons instead of squares.
\end{exlist}
\end{exercisecols}

%========================================================= 11-5
\section*{11--5\quad LIMITS}
\addcontentsline{toc}{section}{11--5\quad Limits}

To investigate circumference and area of circles, it is necessary for us to
discuss \emph{limits}. Precise definitions relating to limits are difficult to
grasp, and a careful study is first made in college calculus courses. For our
purposes, only a little information about limits is required. All our
statements about limits are statements about \emph{numbers}. Therefore you may
regard these statements as giving information about numbers that we will
\emph{assume} or postulate.

\begin{definition}
Suppose that for each positive integer $n$, there is associated in some way, a
number $x_n$. The numbers $x_1$, $x_2,\ldots, x_n,\ldots$ then form an infinite
sequence. The number $x_n$ is called the $n$th term of the sequence. We shall
indicate the sequence by $\{x_n\}$.
\end{definition}

For example, let the $n$th term of a sequence be given by $x_n = 2n$. The
positive integer 1 gives the first term $2 \cdot 1$. The positive integer 2
gives the second term $2 \cdot 2$. The third term is $2 \cdot 3$, and so on.
The sequence given by $x_n = 2n$, then, would look like this:

\begin{center}
\begin{tabular}{r@{\hspace{1.2em}}ccccccc}
Positive integers $\rightarrow$ & 1 & 2 & 3 & 4 & $\ldots$ & $n$ & $\ldots$\\[1pt]
 & $\updownarrow$ & $\updownarrow$ & $\updownarrow$ & $\updownarrow$ & &
   $\updownarrow$ & \\[1pt]
Terms of the sequence $\rightarrow$ & 2 & 4 & 6 & 8 & $\ldots$ & $2n$ &
   $\ldots$\,.
\end{tabular}
\end{center}

\exercisegroup{11--7}
\begin{exercisecols}
\begin{exlist}
\item Write the first six terms of the sequence whose $n$th term is given by
      $x_n = 1/n$; by $x_n = \frac{1}{2}n$; by $x_n = \frac{1}{5}n$.
\item Write the 100th term for each of the sequences in Exercise 1. When $n$ is
      very large, what number are the terms close to?
\item Write the first 10 terms of the sequence whose $n$th term is
      $x_n = (n+1)/n$. Write the 100th term; the 1000th. What number are these
      terms close to when $n$ is extremely large?
\item A sequence is given by $x_n = 2/(n+3)$. Write the first six terms; the
      100th term; the millionth term. Is there a number that the terms of the
      sequence are close to for large values of $n$?
\item Consider a regular polygon of $n$ sides inscribed in a circle. Let $p_n$
      be the perimeter of the polygon. Does this define a sequence?
\item A sequence is given by $x_n = 2/\sqrt{n}$. Write the first six terms. Is
      there a number that the terms of the sequence are close to for large
      values of $n$?
\item Imagine $\sqrt{2}$ written as an infinite decimal,
      $\sqrt{2} = 1.4142\ldots$. Define a sequence $x_n$ as follows: $x_n =$
      the decimal obtained by cutting off the decimal expansion of $\sqrt{2}$
      after the $n$th decimal place. Then $x_1 = 1.4$, $x_2 = 1.41$, etc. What
      is $x_4$? $x_5$? To what number is $x_n$ close when $n$ is large?
\end{exlist}
\end{exercisecols}

We can now define the \emph{limit} of a sequence, an idea that should be clear
from the above exercises. We shall be content with expressing this idea in an
intuitive way, but it should be mentioned that a more complicated definition is
necessary in order to \emph{prove} theorems about limits. We \emph{assume} that
the theorems we need concerning limits are true.

\begin{definition}
The number $L$ is called the \emph{limit} of the sequence $\{x_n\}$ if, for
large values of $n$, the term $x_n$ is necessarily arbitrarily close to $L$. We
then say that $L$ is \emph{the limit of the sequence of terms $x_n$} and write
``$L = \lim\,\{x_n\}$.''
\end{definition}

\exercisegroup{11--8}
In each of the following problems, the $n$th term of a sequence is given. You
are to decide by inspection what the limit of the sequence is. Note that a
sequence does not have to possess a limit.
\begin{exercisecols}
\begin{exlist}
\item $x_n = 1/n$.
\item $x_n = 1/10n$.
\item $x_n = (n+1)/n$.
\item $x_n = \sqrt{n}$.
\item $x_n = 5$.
\item $x_n = 2n/(n+1)$.
\item $x_n = n + 2$.
\item[\stex 8.] $x_n = (2n + 100)/(4n + 50)$.
\item[\stex 9.] $x_n = (n + 1)/n^2$.
\item[\stex 10.] $x_n = (n^2 + 5n + 6)/[2n(n + 2)]$.
\end{exlist}
\end{exercisecols}

When we discuss the circumference and area of circles, we will need the
following two theorems on limits. We assume that these theorems are true, and
thus treat them as postulates. But they are not postulates about our
\emph{geometry}; they are postulates about the \emph{number system}.

\begin{theorem}
Consider two sequences with $n$th terms $x_n$ and $y_n$, respectively, and
suppose that $\lim\,\{x_n\} = X$ and $\lim\,\{y_n\} = Y$. Then the sequence
whose $n$th term is $x_n y_n$ has a limit, and $\lim\,\{x_n y_n\} = XY$.
\end{theorem}

\begin{theorem}
With the same notation as in Theorem 11--10, and the additional hypothesis that
$y_n$ and $Y$ are not zero, then the sequence whose $n$th term is $x_n/y_n$ has
a limit, and $\lim\,\{x_n/y_n\} = X/Y$.
\end{theorem}

\exercisegroup{11--9}
\begin{exercisecols}
\begin{exlist}
\item Verify Theorem 11--10 if $x_n = (2n + 1)/n$ and $y_n = n/(n + 1)$.
\item If $\lim\,\{x_n\} = 2$, what is $\lim\,\{(x_n)^2\}$?
\item If $\lim\,\{x_n\} = 3$ and $\lim\,\{y_n\} = 4$, what is
      $\lim\,\{x_n/y_n\}$? $\lim\,\{x_n y_n\}$? $\lim\,\{y_n/x_n\}$?
\item If $x_n$ is ``very close'' to 5 and $y_n$ is ``very close'' to 7, what is
      $x_n y_n$ close to? $x_n/y_n$? $y_n/x_n$?
\end{exlist}
\end{exercisecols}

%========================================================= 11-6
\section*{11--6\quad CIRCUMFERENCE}
\addcontentsline{toc}{section}{11--6\quad Circumference}

From your everyday experience, you already have an idea of the perimeter or
circumference of a circle. For example, if you were asked to find the
circumference of a wheel, you might wrap a string around the wheel, then pull
it out straight, and measure its length. But such a practical method is not
adequate for geometry, since we cannot describe this process in our geometry.
It is necessary that we first \emph{define} circumference. The definition must
satisfy our intuition about what we \emph{feel} the circumference should be,
and should be one that can readily be \emph{used} to prove theorems.

Observe that we had no difficulty in defining the perimeter of a polygon,
because the polygon consisted of a finite number of segments, each of which had
a length. A circle contains no segment of a line, so we cannot define its
perimeter so simply. To arrive at our definition, we need the following
theorem, the truth of which we assume.

\begin{theorem}
For a given circle, let $p_n$ be the perimeter of a regular polygon of $n$
sides inscribed in the circle. The numbers $p_n$ $(n = 3, 4, 5,\ldots)$ form a
sequence that has a limit.
\end{theorem}

\begin{XIremarks}
The proof of this theorem involves two steps. First, it is shown that the
sequence is increasing, that is, each term is less than the one that follows
it. Then it is shown that the sequence is \emph{bounded}, that is, the terms do
not become indefinitely large. It then follows, as a fundamental property of
real numbers, that such a sequence has a limit.
\end{XIremarks}

It is natural to ask how we know that regular polygons with $n$ sides exist for
all positive integers $n \XIgeqq 3$. We do know that there are regular polygons
with the number of sides equal to $2^2$, $2^3$, $2^4$, etc., for we can
inscribe a square in the circle, and, by bisecting the central angles, we get a
regular octagon, and then a regular 16-gon, and so on, as in Fig.~11--10.

\FIGXITEN

Therefore, the perimeters of these special regular polygons form a sequence
$p_4$, $p_8$, $p_{16}$, $p_{32},\ldots$, and by Theorem 11--12 and the
definition of limit, this sequence has a limit.

We can now give a definition of circumference. Observe that in order to
\emph{define} the circumference of a circle, we are \emph{necessarily forced}
to use limits.

\begin{definition}
The \emph{circumference $C$} of a circle is the limit of the sequence of
perimeters $p_n$ of regular inscribed polygons. $C = \lim\,\{p_n\}$.
\end{definition}

\exercisegroup{11--10}
\begin{exercisecols}
\begin{exlist}
\item For a given circle, show that the perimeter of an inscribed square is
      less than that of an inscribed octagon.
\item Find the perimeter of a square inscribed in a circle of radius 1. Hence
      show that the circumference of the circle is more than $4\sqrt{2}$.
\item By inscribing a regular hexagon in a circle of radius 1, show that the
      circumference of the circle is more than 6.
\end{exlist}
\end{exercisecols}

The fundamental result on the circumferences of circles is the following:

\begin{theorem}
The circumferences of two circles are proportional to their radii (or to their
diameters). That is, if the circles have circumferences $C$ and $C'$, radii $r$
and $r'$, and diameters $d$ and $d'$, then $C/C' = r/r' = d/d'$.
\end{theorem}

\noindent\emph{Proof.}
\begin{steps}
\step{$C/C' = \lim\,\{p_n\}/\lim\,\{p'_n\}$, where $p_n$, $p'_n$, are the
      perimeters of the regular inscribed polygons of $n$ sides in the
      circles.}{Definition 11--7.}
\step{$C/C' = \lim\,\{p_n/p'_n\}$.}{Theorem 11--11 on limits.}
\step{$p_n/p'_n = r/r' = d/d'$.}{Theorem 11--7.}
\step{$\lim\,\{p_n/p'_n\} = \lim\,\{r/r'\}$.}{$p_n/p'_n = r/r'$ for all $n$.}
\step{$\lim\,\{r/r'\} = r/r'$.}{$r$ and $r'$ are the same for all $n$.}
\step{$\lim\,\{p_n/p'_n\} = r/r' = d/d'$.}{Statements 3, 4, and 5.}
\step{$C/C' = r/r' = d/d'$.}{Statements 2 and 6.}
\end{steps}

\begin{remark}
Observe that this theorem about circumferences is basically a theorem about
similar polygons.
\end{remark}

\begin{theorem}
If $C$ is the circumference of a circle, and $d$ is its diameter, the ratio
$C/d$ is the same for all circles.
\end{theorem}

\noindent\emph{Proof.} This theorem is really just a corollary to Theorem
11--13. But it is so important that it is given a place by itself. The proof is
as follows. From Theorem 11--13, we have $C/C' = d/d'$. This proportion implies
the proportion $C/d = C'/d'$. This last equality says that $C/d$ is the same
number for any two circles. This completes the proof.

\begin{definition}
The number $C/d$, which is the same for all circles, will be denoted by the
Greek letter $\pi$ (pi).
\end{definition}

\begin{corollary}
$C = \pi d = 2\pi r$.
\end{corollary}

\begin{remark}
The definition of $\pi$ does not immediately give its numerical
value. So a problem remains---namely that of calculating $\pi$ to any desired
degree of accuracy. We examine later a bit of the history of this number and
methods for calculating its value. At the moment, and for purposes of the
exercises, we mention two approximate values commonly used for $\pi$, namely
22/7 and 3.1416.
\end{remark}

\exercisegroup{11--11}
\begin{exercisecols}
\begin{exlist}
\item Find, approximately, the circumference of a circle whose radius is 7 in.;
      2.131 in.; 49 ft.
\item Find, approximately, the radius of a circle whose circumference is 66 ft;
      12.125 in.; 100.00 ft.
\item Give an exact expression for the area of a square whose perimeter is
      equal to the circumference of a circle of radius 1.
\item Suppose that a circular steel band surrounds the earth at the equator
      (radius $=$ 4000 mi). If the band is made longer by inserting a piece
      1 ft long, about how far will the band be from its previous position?
\item A tire on a car has a diameter of 26 in. If it revolves at 10 revolutions
      per second, what is the approximate speed of the car in miles per hour?
\item About how many revolutions per minute does a car wheel make (diameter
      $=$ 28 in.) when the car is traveling at 60 mph?
\item The radius of one circle is three times the radius of a second. How do
      their circumferences compare?
\end{exlist}
\end{exercisecols}

%========================================================= 11-7
\section*{11--7\quad COMPUTATION OF PI}
\addcontentsline{toc}{section}{11--7\quad Computation of $\pi$}

There are many ways of computing $\pi$, some much more convenient than others.
The obvious way is to compute the perimeters $p_n$ of regular polygons with a
large number of sides. Then $\pi$ would be given approximately by $p_n/d$. This
is a rather tedious process, and we use it here to get a very rough estimate of
$\pi$.

\FIGXIELEVEN

Consider a regular inscribed hexagon in a circle of radius 1, as shown in
Fig.~11--11. The perimeter of the hexagon is 6. Increasing the number of sides
will increase the perimeter. Hence $C > 6$, so that $\pi = C/d > 6/2 = 3$, and
$\pi > 3$.

If we consider a regular circumscribed hexagon, we find that its perimeter is
$4\sqrt{3}$, and since $C$ is less than this perimeter,
\[
\pi = \frac{C}{d} < \frac{4\sqrt{3}}{2} = 2\sqrt{3} < 3.465,
\]
and $\pi < 3.465$.

In this way, we have obtained a very rough value for $\pi$. It lies somewhere
between 3 and 3.5. By these same methods, Archimedes (3rd century B.C.) showed
that $\pi$ is between 3 1/7 and 3 10/71. The Babylonians used 3 as the value of
$\pi$, and there is a verse in the Bible that implies that the writer believed
$\pi = 3$.

By far the most convenient and rapid methods for calculating $\pi$ are based on
formulas derived in calculus.\XIsymfoot{A formula from calculus that can be
used to calculate $\pi$ is
\[
\pi = 4\left(1 - \tfrac{1}{3} + \tfrac{1}{5} - \tfrac{1}{7} + \tfrac{1}{9}
- \tfrac{1}{11} + \cdots\right).
\]}
We cannot take up these formulas here, but it is interesting to know that in
the 17th century a Scot by the name of Shanks,\XIsymfoot{One of your authors
makes no claim to be a descendant of this misguided computer.} using one of
these formulas, computed $\pi$ to 707 decimal places. His reasons for doing
this are at best obscure. More recently, a group working on one of the large
electronic computers set up a program for the calculation of $\pi$ on the
computer and obtained, in a few hours, $\pi$ to more than 2000 decimal places!
This demonstrates the extreme speed of modern computing machines. The
corresponding job by hand would take hundreds of years. The correct value of
$\pi$ to 14 decimal places is 3.14159265358979.

The number $\pi$ is an irrational number. A proof of this is not easy, and was
given first by the German mathematician Johann Lambert in 1766.

\exercisegroup{11--12}
\begin{exercisecols}
\begin{exlist}
\item By considering inscribed and circumscribed squares, show that
      $2\sqrt{2} < \pi < 4$.
\item[\stex 2.] By considering an inscribed regular octagon, show that
      $\pi > 3.06$.
\item[\stex 3.] Find the perimeter of a regular 12-gon in a circle of radius 1.
      Hence find an approximate value for $\pi$. Can you continue to find the
      perimeter of a regular 24-gon?
\item[4.] What does the irrationality of $\pi$ imply about its decimal
      representation?
\end{exlist}
\end{exercisecols}

%========================================================= 11-8
\section*{11--8\quad AREA OF CIRCLES}
\addcontentsline{toc}{section}{11--8\quad Area of circles}

In studying area we first obtained the area of a rectangle by counting squares.
Areas of triangles may be obtained by forming parallelograms and observing that
the triangle is half the area of the parallelograms. Areas of polygons were
obtained by cutting up the polygons into triangles. To define what we mean by
the area of a circle, it is necessary to use again the idea of limit, just as it
was necessary to use limits in order to define circumference. We all have a
fairly clear idea of what we \emph{feel} the area of a circle should be, and we
use this impression to define area. Since the area of the circle should be
close to the area of the inscribed polygon with a large number of sides, we
first need to know that the limit of the areas of inscribed polygons exists.
This fact is explicitly stated as a theorem below, which we accept without
proof.

\begin{theorem}
Given a circle $C$, let $S_n$ denote the area of an inscribed regular polygon
of $n$ sides; then $\lim\,\{S_n\}$ exists. In other words, there is a number
$S$, such that the term $S_n$ is arbitrarily close to $S$ for large $n$.
\end{theorem}

\begin{XIremarkn}{1}
We know, from Section 11--6, that regular polygons exist with the number of
sides equal to powers of 2. In using Theorem 11--15, we can consider only
polygons that we know to exist.
\end{XIremarkn}

\begin{XIremarkn}{2}
Theorem 11--15 is a theorem about numbers. It has little geometric content,
so we need not feel that it is another geometric postulate. If we consider only
polygons with 4, 8, 16, 32,\ldots\ sides, it is easy to show that each area in
our sequence $S_4$, $S_8$, $S_{16}$, $S_{32},\ldots$ is less than the one that
follows. On the other hand, we can see that each of these polygons has an area
less than the area of a circumscribed square. This indicates that our sequence
is not only increasing, but is ``bounded above.'' It is a theorem about numbers
that such a sequence has a limit.
\end{XIremarkn}

\begin{definition}
The area $S$ of a circle is the limit of the areas of inscribed regular
polygons. That is, $S = \lim\,\{S_n\}$.
\end{definition}

We can now prove the main theorem about the area of a circle. But because the
definition of area uses the idea of limit, we must use some theorems about
limits. In particular, we need to know the following fact about the apothems of
the inscribed polygons.

\begin{theorem}
Let $a_n$ be the apothem of an inscribed regular polygon of $n$ sides in a
circle of radius $r$, then $\lim\,\{a_n\} = r$.
\end{theorem}

\FIGXITWELVE

\noindent\emph{Proof.} In Fig.~11--12, $\sg{AC}$ is the length of a side of the
polygon, and is a small number when the number of sides of the polygon is
large. By the Pythagorean theorem, $a_n = \sqrt{r^2 - \sg{AB}^2}$. Thus, when
$n$ is very large, $\sg{AB}^2$ is very small, and $r^2 - \sg{AB}^2$ is very
near $r^2$. Hence $a_n$ will be near $r$. This is exactly what we mean when we
say that $\lim\,\{a_n\} = r$.

\begin{theorem}
The area of a circle is equal to one-half the circumference times the radius.
$S = \frac{1}{2}$ (circumference) $\times$ (radius).
\end{theorem}

\noindent\emph{Proof.}
\begin{steps}
\step{$S = \lim\,\{S_n\}$, where $S_n$ is the area of the regular inscribed
      polygon of $n$ sides.}{Definition.}
\step{$S_n = \frac{1}{2}p_n \cdot a_n$, where $p_n$ is the perimeter, and $a_n$
      is the apothem, of the polygon.}{Theorem 11--8.}
\step{$\lim\,\{S_n\} = \frac{1}{2}(\lim\,\{p_n\}) \cdot
      (\lim\,\{a_n\})$.}{Statement 2 and Theorem 11--10 on limits.}
\step{$\lim\,\{p_n\} = C =$ circumference.}{Definition of circumference.}
\step{$\lim\,\{a_n\} = r =$ radius.}{Theorem 11--16.}
\step{$\lim\,\{S_n\} = \frac{1}{2}C \cdot r$.}{Statements 3, 4, and 5.}
\step{$S = \frac{1}{2}C \cdot r$.}{Statements 1 and 6.}
\end{steps}

\begin{corollary}
The area of a circle of radius $r$ is $\pi r^2$.
\end{corollary}

\begin{corollary}
The areas of two circles are proportional to the squares of their radii (or
diameters).
\end{corollary}

\noindent The proofs are left to the student.

% multicols dropped: this group carries Figs 11-13, 11-14, 11-15 and 11-16.
% They were already sitting beside the exercises that name them, but inside a
% 52 mm column, which printed them smaller than the book's own; each is now an
% \exfig on the full measure.
\exercisegroup{11--13}
\begin{exercisecols}
\begin{exlist}
\item Two circles have radii of 5 in. and 2.5 in. Find their areas. What is the
      ratio of their areas?
\item An annular ring (Fig.~11--13) has inner radius 2 ft and outer radius
      4 ft. Find its area.
\exfig{\FIGXITHIRTEEN}
\item Figure 11--14 is composed of a circle and semicircles. If
      $\sg{AB} = \sg{BC} = 2$, $B$ is the center of the circle, and $AC$ is a
      diameter, find the area of the shaded portion.
\exfig{\FIGXIFOURTEEN}
\item Show that the area of the inscribed circle is half the area of the
      circumscribed circle of the square (Fig.~11--15).
\exfig{\FIGXIFIFTEEN}
\item The area of a circle is 154 in\textsuperscript{2}. Find its radius.
      (Use $\pi = 22/7$.)
\item The area of a circle is 100 ft\textsuperscript{2}. Find its radius.
\item Semicircles are constructed on the sides of a right triangle, as in
      Fig.~11--16. Show that the area of the largest is the sum of the areas of
      the other two.
\exfig{\FIGXISIXTEEN}
\item A circle has its circumference equal to its area. What is its radius?
\item A rectangular barn is 30 ft by 40 ft. A goat is tethered at one corner by
      a rope 40 ft long. Over what area can he graze? Over what area can he
      graze if the rope is 70 ft long?
\item The circumference of a circle is 50 ft. What is its area?
\item The area of a circle is 288 in\textsuperscript{2}. What is the area if
      the unit of length is a foot?
\item On Mars the standard unit of length is the \emph{pfutth}, which is 0.88
      as long as our foot. A Martian states that the area of a circle is 16
      square \emph{pfutths}. What is the area in square feet?
\item[\stex 13.] Show that the area of a circle is the limit of the areas of
      circumscribed regular polygons.
\item[14.] The area of a circle is 49/4 times the area of another circle. What
      is the ratio of their radii?
\end{exlist}
\end{exercisecols}

%========================================================= review
\section*{REVIEW OF CHAPTER 11}
\addcontentsline{toc}{section}{Review of Chapter 11}

The following definitions were given in this chapter. Be sure that you know
them.

\begin{center}
\begin{minipage}{0.66\linewidth}
inscribed polygon (circle)\\
circumscribed polygon (circle)\\
tangent\\
center, radius, apothem of regular polygons\\
sequence, limit\\
circumference\\
$\pi$\\
area of a circle
\end{minipage}
\end{center}

The following is a list of the more important theorems of the chapter. Give
complete statements of these theorems. Then try to give a proof for each. Do
not refer to the text until you need help.

\begin{exlist}
\item A regular polygon has a circumscribed circle.
\item A line tangent to a circle is perpendicular to a radius.
\item Converse of 2 (almost).
\item Regular polygons are similar if and only if they have the same number of
      sides.
\item The perimeters of similar polygons are proportional to their radii.
\item The area of a regular polygon is one half the product of apothem and
      perimeter.
\item The circumferences of two circles are proportional to their radii.
\item The area of a circle is equal to one half the circumference times the
      radius.
\end{exlist}

\begin{center}\bfseries Miscellaneous Questions\end{center}
\nopagebreak

\begin{exlist}
\item Why is it important in this chapter to know that regular polygons of the
      same number of sides are similar?
\item Why is it important that circumference is proportional to the radius?
      (Tell why this enables us to define $\pi$.) This theorem depends on what
      theorems on similar polygons?
\item Tell why we must \emph{define} circumference and area for circles.
\item In proving that area $= \pi r^2$, what theorems about polygons did we
      need? what theorems on limits?
\end{exlist}

\end{document}
